\section{From Methodology to Measurement}
\label{sec:transition_dnds}

The CNN-based framework described in the preceding sections was designed to overcome the limitations of existing analytical methods while retaining their strengths.
Recovering the sub-threshold $dN/dS$ quickly and with calibrated uncertainties opens several scientific directions.

First, a robust measurement of the source-count distribution constrains the faint-end population of astrophysical emitters --- blazars, star-forming galaxies, and millisecond pulsars --- that collectively generate the bulk of the unresolved gamma-ray background.
Quantifying their contribution is essential before any residual can be attributed to non-standard sources.
Second, the recovered $dN/dS$ provides the empirical prior for the probabilistic source cataloging framework developed in Chapter~\ref{ch:7}: the flux distribution of sub-threshold sources directly determines the posterior probability that any given sky pixel contains a faint, as-yet-unresolved source.
Third, the methodology itself is readily extensible.
An energy-dependent $dN/dS$ reconstruction across multiple energy bands would enable spectral decomposition of the unresolved sky into individual source classes --- and, in principle, a search for spectral signatures inconsistent with known astrophysical populations.
This extension, while beyond the scope of the present work, represents a natural next step toward the dark matter cross-correlation searches pursued in Chapter~\ref{ch:8}.

% The analysis presented below should therefore be understood as a proof of principle for the deep-learning approach to population-level gamma-ray inference.
% We apply the trained CNN to 14~years of \textit{Fermi}-LAT data in the 1--10~GeV energy band and recover the source-count distribution down to fluxes approximately 50 times below the catalog detection threshold, demonstrating that the method produces stable, well-calibrated results across a wide dynamic range.
\blue{Applied to real \textit{Fermi}-LAT data, the framework recovers the source-count distribution down to fluxes roughly 50 times below the catalog detection threshold, with stable, well-calibrated results across a wide dynamic range.}
\blue{The recovered distribution is broadly consistent with independent determinations in the literature. These range from the 1pPDF and Compound Poisson Generator analyses of Zechlin, Cuoco et al.~\cite{Cuoco-1pdf} to the recent simulation-based inference reconstruction of Eckner et al.~\cite{Eckner:2025waa}, which probe the same high-latitude sky with different SBI techniques.}
\blue{The remainder of this chapter is based on \cite{Amerio:2023uet} and presents the analysis in full.}
